\documentclass{article}

\usepackage[submission]{colm2026_conference}
\usepackage{booktabs}
\usepackage{graphicx}
\usepackage{microtype}
\usepackage{hyperref}
\usepackage{url}
\usepackage{xcolor}
\usepackage{lineno}
\usepackage{amsmath}

\definecolor{darkblue}{rgb}{0,0,0.5}
\hypersetup{colorlinks=true, citecolor=darkblue, linkcolor=darkblue, urlcolor=darkblue}

\title{Pseudonymized but Linkable:\\
Cross-Document Re-identification Risks in LLM-Ready Sensitive Text Corpora}

\author{}
\date{}

\newcommand{\method}{LinkGuard}
\newcommand{\bench}{CrossDoc-PrivacyBench}

\begin{document}
\ifcolmsubmission
\linenumbers
\fi
\maketitle

\begin{abstract}
Sensitive text corpora are commonly transformed one document at a time before retrieval, evaluation, sharing, or model adaptation.
This workflow removes obvious direct identifiers, but it can leave a corpus-level privacy failure mode: records with no names, emails, or account numbers may remain linkable through repeated quasi-identifiers such as role, region, schedule, institution type, family context, health context, financial context, and rare events.
We introduce \bench, a fully synthetic benchmark with 120 controlled personas, 480 documents, risk tiers, utility labels, and decoy auxiliary profiles.
We evaluate document-local transformations under linkage, auxiliary-profile matching, RAG-style exposure, quasi-identifier recovery, and utility metrics.
Direct redaction leaves auxiliary-profile matching unchanged at 0.71 top-1, while a Presidio-style PII baseline remains highly matchable at 0.59.
We propose \method, a simple corpus-aware generalization heuristic that estimates cross-document uniqueness and generalizes high-risk quasi-identifiers until a target anonymity set is reached.
At target $k=5$, \method{} reduces auxiliary top-1 to 0.04 while preserving issue classification and retrieval Recall@5 at 1.00; aggressive redaction reaches similar privacy but drops issue accuracy to 0.35 and retrieval Recall@5 to 0.31.
GPT-5.5 auxiliary and RAG-generation stress audits preserve this ordering, reinforcing that LLM-ready data transformations should be evaluated for corpus-level linkage resistance, not only document-local PII removal.
\end{abstract}

\section{Introduction}

Sensitive data is often most useful when it is not isolated.
A hospital may have clinical messages, billing notes, scheduling records, and accommodation requests about the same patient.
A legal-aid organization may have intake notes, benefits appeals, landlord disputes, and follow-up emails about the same client.
A company may have support tickets, HR notes, finance requests, and internal chat snippets about the same employee or customer.
Foundation-model workflows, retrieval-augmented generation (RAG), evaluation-set construction, model adaptation, analytics, and data sharing, therefore often operate over \emph{multi-document} corpora rather than one-off documents.

The dominant operational response is document-local de-identification.
A pipeline removes names, emails, phone numbers, mailing addresses, account numbers, and sometimes other spans that look individually identifying.
This is necessary, but it is not sufficient.
A rare combination such as a region, occupation, institution type, care schedule, family constraint, legal context, and local event may be harmless in one document yet become a stable signature when repeated across documents.
Pseudonymization can make this worse: a stable pseudonym is useful for continuity, but it is also a linkage handle.
Per-document pseudonyms remove the explicit handle, but they do not remove the repeated quasi-identifier structure.

This paper studies a focused question:
\emph{after document-local anonymization or pseudonymization, can an attacker still link multiple transformed documents to the same person and match them to auxiliary background profiles?}
We intentionally avoid real-person data and real re-identification attempts.
Instead, we build a controlled synthetic benchmark in which every persona, document, decoy, and utility label is generated synthetically to give us safe ground truth for controlled privacy auditing.

Our contributions are:
\begin{itemize}
    \setlength{\itemsep}{0pt}
    \setlength{\parsep}{0pt}
    \setlength{\topsep}{2pt}
    \item \textbf{Benchmark:} \bench, a synthetic cross-document privacy benchmark with 120 personas, 480 documents, risk tiers, utility labels, templates, and decoy auxiliary profiles.
    \item \textbf{Audit protocol:} pairwise linkage, fixed-$K$ clustering, auxiliary-profile matching, quasi-identifier recovery, RAG-style exposure, and utility retention.
    \item \textbf{Finding:} direct redaction, Presidio-style PII removal, document-local anonymization, and stable pseudonyms all leave strong residual linkage risk.
    \item \textbf{Defense:} \method, a simple corpus-aware generalization heuristic that improves the synthetic privacy--utility tradeoff over blanket aggressive redaction.
\end{itemize}

\section{Problem Setup and Threat Model}

\paragraph{Setting.}
A data holder has a sensitive text corpus with multiple records per person, household, customer, patient, client, or employee.
The corpus is transformed before it is indexed for RAG, shared internally, released for evaluation, or used in model-adaptation workflows.
We consider transformations that operate on text, including direct redaction, pseudonymization, document-local anonymization, and corpus-aware generalization.

\paragraph{Attacker knowledge.}
The attacker sees the transformed corpus.
For the auxiliary-profile task, the attacker also receives a small candidate set containing one true synthetic auxiliary profile and decoys.
This simulates background knowledge such as an external directory, leaked structured metadata, public biographical context, or another internal table.
The attacker does not receive the original direct identifiers, and the candidate labels are neutralized to avoid leaking synthetic persona IDs.

\paragraph{Attacker goals.}
We evaluate three defensive questions.
First, can an attacker tell which transformed documents belong to the same synthetic persona?
Second, can the attacker match a cluster of transformed documents to the correct auxiliary profile among decoys?
Third, can profile-like queries retrieve target documents from a transformed corpus, creating a RAG exposure channel?

\paragraph{Non-goals and ethics.}
We do not collect, process, or attempt to re-identify real people.
All personas, documents, auxiliary profiles, and decoys are synthetic.
The benchmark is intended for defensive auditing of transformation methods, not for operational re-identification.
Any extension to real data would require authorization, data minimization, security review, and domain-specific ethical oversight.

\section{\bench: A Synthetic Cross-Document Benchmark}

\paragraph{Personas.}
\bench{} contains 120 synthetic personas.
Each persona has direct identifiers: name, email, phone number, address, and account ID, and a structured set of quasi-identifiers: region, city, occupation, employer type, education, family structure, medical context, financial context, legal context, schedule pattern, community affiliation, and rare event.
Each persona also has utility labels for four domains.
We use deterministic generation so that experiments are reproducible and no real individual is represented.

\paragraph{Risk tiers.}
Personas are assigned to three tiers.
T1 documents contain mostly generic quasi-identifiers, so direct PII is the main privacy risk.
T2 documents contain moderately specific but partly generalized context, such as a city within a region, employer category, broad health or legal context, and schedule constraints.
T3 documents contain high-linkage combinations, such as exact city, occupation, employer type, education, family structure, treatment or schedule context, affiliation, and rare event.
This tiering lets us check whether an attack and defense are sensitive to increasing quasi-identifier specificity.

\paragraph{Documents.}
Each persona has four documents spanning healthcare, legal, financial support, and HR.
The documents are short service or support records with a direct-contact line and a task-oriented body.
Each document includes an issue label, such as medication refill, benefits appeal, payment deferral, housing issue, accommodation request, or schedule change.
The benchmark therefore contains 480 documents total and 384 held-out test documents from 96 held-out personas.

\paragraph{Auxiliary profiles and decoys.}
For each persona, we construct a ten-candidate auxiliary set.
One candidate is the true profile; nine are decoys selected to share some attributes with the target, such as region, occupation family, or legal context.
This makes matching harder than a trivial exact lookup and lets us report top-1, top-3, and mean reciprocal rank.
The auxiliary-profile task is the central privacy test because it captures both linkage and background-knowledge use.


\section{Transformations}

We compare eight conditions.
C0 is the original synthetic corpus.
C1 performs direct redaction of names, first and last names, email addresses, phone numbers, mailing addresses, and account IDs.
C1b applies a Presidio-based PII analyzer plus local regex cleanup for the synthetic account and address formats.
C2 uses stable pseudonyms across the corpus, replacing direct identifiers and some locations/institutions with consistent tokens.
C3 resets pseudonym mappings per document.
C4 is a document-local anonymization proxy that removes direct identifiers and generalizes selected local spans, such as city and education, without seeing neighboring documents.
C5 is \method.
C6 is aggressive redaction, which suppresses most potentially identifying context, including quasi-identifiers and task issue phrases.

\subsection{\method: Corpus-Aware Generalization}

Document-local anonymizers ask whether a span is identifying inside one document.
\method{} asks whether repeated details make a persona distinguishable in the corpus.
It operates over persona groups in the synthetic benchmark, but the same logic could be applied to known customer, patient, case, or account groupings in a controlled deployment.

For each persona and quasi-identifier field, \method{} defines a three-level transformation ladder:
\[
\text{exact value} \rightarrow \text{semantic category} \rightarrow \text{generic suppression}.
\]
For example, a city can become a region and then ``a broad location''; an occupation can become a job family and then ``general worker''; a rare event can become a local-event category and then ``a local event.''
The method starts from direct redaction and estimates how many personas share the current generalized representation.
If the estimated anonymity set is smaller than the target $k$, \method{} selects another field to generalize.
The selection score balances rarity, a specificity prior, and an approximate utility cost:
\[
s(f) =
\frac{\mathrm{rarity}(f) + \mathrm{specificity}(f) + 1}
{\max(\mathrm{utility\_cost}(f), 0.1)}.
\]
Fields with rarer values, higher specificity, and lower expected utility cost are generalized first.
The process stops when the estimated anonymity set reaches target $k$ or all fields have been generalized or suppressed.
This is a heuristic risk-budgeting procedure, the estimate depends on the modeled population, the field schema, and the attacker's auxiliary knowledge.

\section{Attacks and Metrics}

\paragraph{Pairwise linkage.}
We compute signature features over transformed documents, including explicit pseudonym handles and quasi-identifier phrases.
Pairwise linkage is evaluated as same-person binary classification over held-out document pairs.
We select a threshold on validation personas and report test F1.
This metric is especially useful for showing that stable pseudonyms can become linkage handles.

\paragraph{Fixed-$K$ clustering.}
We also report fixed-$K$ clustering adjusted Rand index, where $K$ is the true number of held-out personas.
Threshold graph clustering is brittle outside stable-handle settings, so we treat pairwise linkage and auxiliary matching as the main privacy evidence.

\paragraph{Auxiliary-profile matching.}
For each held-out persona, the attacker ranks the ten auxiliary candidates using the transformed documents.
The main local attacker uses word TF-IDF similarity; robustness checks use character TF-IDF, hybrid TF-IDF, and a field-weighted matcher that explicitly scans exact and generalized schema fields.
We report top-1, top-3, and mean reciprocal rank.
Lower values indicate better privacy.

\paragraph{Exact quasi-identifier recovery.}
Auxiliary matching gives an end-to-end privacy signal, but it can hide which information leaked.
We therefore scan transformed documents for exact recovery of quasi-identifier fields such as city, occupation, employer type, education, family context, medical context, financial context, legal context, schedule pattern, affiliation, and rare event.
This is a structured profile-reconstruction diagnostic.

\paragraph{RAG exposure.}
To model retrieval exposure, we issue synthetic profile queries against the transformed corpus and measure whether target-persona documents appear in the top retrieved results.
We report Hit@5, whether multiple target documents appear in the top ten, document recall at ten, and reciprocal rank.
This diagnostic asks whether transformed records remain reachable from profile-like background knowledge even when they contain no direct identifiers.

\paragraph{Utility.}
We evaluate lightweight task utility with issue classification, task-oriented retrieval Recall@5, and fact preservation.
The classifiers and retrievers are local TF-IDF models trained on original synthetic documents.
These metrics test whether service-triage information survives transformation; they do not claim to cover all downstream uses.

\section{Results}

\subsection{Main Privacy--Utility Sweep}

Table~\ref{tab:main} summarizes the main local sweep over 96 held-out personas.
Direct redaction leaves auxiliary matching unchanged at 0.7 top-1, confirming that direct PII removal alone does not remove repeated quasi-identifier risk.
The Presidio-style baseline improves over direct redaction but remains highly matchable at 0.6.
The document-local proxy still leaves auxiliary top-1 at 0.6 and exact quasi-identifier recovery at 0.3.
Consistent pseudonymization makes linkage nearly trivial, with pairwise F1 of 0.98, because the stable token itself becomes a cross-document handle.

\begin{table}[t]
\centering
\scriptsize
\setlength{\tabcolsep}{3.2pt}
\begin{tabular}{lrrrrrrr}
\toprule
Cond. & Pair & Aux@1 & Exact & Issue & Ret@5 & Facts & Edit \\
\midrule
C0 Orig & 0.96 & 0.71 & 0.39 & 1.00 & 1.00 & 0.68 & 0.00 \\
C1 Redact & 0.21 & 0.71 & 0.39 & 1.00 & 1.00 & 0.68 & 0.25 \\
C1b Presidio & 0.14 & 0.59 & 0.33 & 1.00 & 1.00 & 0.68 & 0.35 \\
C2 Stable & 0.98 & 0.43 & 0.27 & 1.00 & 1.00 & 0.68 & 0.30 \\
C3 PerDoc & 0.15 & 0.43 & 0.27 & 1.00 & 1.00 & 0.68 & 0.31 \\
C4 Local & 0.16 & 0.60 & 0.30 & 1.00 & 1.00 & 0.68 & 0.29 \\
C5 \method{} & 0.02 & 0.04 & 0.00 & 1.00 & 1.00 & 0.70 & 0.54 \\
C6 Agg & 0.02 & 0.05 & 0.00 & 0.35 & 0.31 & 0.39 & 0.68 \\
\bottomrule
\end{tabular}
\caption{Main held-out results. Pair is pairwise linkage F1; Aux@1 is auxiliary-profile top-1 accuracy; Exact is exact quasi-identifier recovery; Issue and Ret@5 are utility metrics; Facts is lightweight fact preservation; Edit is character-level edit ratio. Lower Pair, Aux@1, Exact, and Edit are better for privacy or fidelity; higher Issue, Ret@5, and Facts are better for utility.}
\label{tab:main}
\end{table}

\method{} reduces auxiliary top-1 to 0.04 and exact recovery to 0.0 while preserving issue accuracy and retrieval Recall@5 at 1.0.
Aggressive redaction reaches a similar auxiliary top-1 of 0.05, but issue accuracy drops to 0.35 and retrieval Recall@5 to 0.31.
Thus, in this synthetic sprint, corpus-aware generalization occupies a better privacy--utility point than blanket suppression.
A stricter body-only utility stress test tells the same story: \method{} scores 0.88, direct redaction 0.87, and aggressive redaction 0.44.

\subsection{Risk Tiers and Signal Ablations}

The benchmark tiers behave as intended.
Under direct redaction, Aux@1 rises from 0.16 on T1 to 0.97 on T2 and 1.00 on T3.
\method{} lowers the corresponding rates to 0.03, 0.00, and 0.09, and reduces T3 exact quasi-identifier recovery to 0.00.
The result indicates that direct redaction is adequate only in the tier where quasi-identifiers were intentionally generic; it fails as soon as the corpus contains repeated, moderately specific context.

Single-signal ablations of direct-redacted documents indicate that role and institution are the largest local drivers.
Removing role and organization drops Aux@1 from 0.71 to 0.50, while removing location drops it to 0.52.
Family, health, rare event, and schedule ablations have little effect individually in this synthetic corpus, but they still matter in combination and under the field-weighted stress attacker.
This is consistent with the central claim: linkage risk is often combinatorial, not span-local.

\begin{table}[t]
\centering
\small
\setlength{\tabcolsep}{5pt}
\begin{tabular}{lrrrr}
\toprule
Cond. & T1 Aux@1 & T2 Aux@1 & T3 Aux@1 & T3 Exact \\
\midrule
C1 Redact & 0.16 & 0.97 & 1.00 & 1.00 \\
C1b Presidio & 0.09 & 0.69 & 1.00 & 0.88 \\
C4 Local & 0.16 & 0.66 & 1.00 & 0.82 \\
C5 \method{} & 0.03 & 0.00 & 0.09 & 0.00 \\
C6 Agg & 0.06 & 0.00 & 0.09 & 0.00 \\
\bottomrule
\end{tabular}
\caption{Risk-tier stratification. Direct and document-local baselines fail most strongly on high-linkage T3 personas; \method{} reduces both auxiliary matching and exact quasi-identifier recovery.}
\label{tab:tier}
\end{table}

\subsection{Sensitivity to the Linkage Risk Budget}

Table~\ref{tab:sensitivity} varies the target anonymity set size used by \method{}.
Moving from direct redaction ($k=1$) to $k=2$ drops Aux@1 from 0.71 to 0.06 and exact recovery from 0.39 to 0.01; $k=5$ removes exact recovery and lowers the field-weighted attacker to 0.24.
Larger $k$ mainly helps the field-weighted stress attacker, which falls to 0.10 at $k=20$, so corpus-aware generalization acts as a tunable risk budget.
A corpus-awareness ablation confirms that true corpus statistics matter: target-$k$ coverage is 1.00 with true planning but 0.92 under shuffled quasi-identifier co-occurrences, while a global level-1 rule leaves field-aware top-1 at 0.98.

\begin{table}[t]
\centering
\scriptsize
\setlength{\tabcolsep}{3pt}
\begin{tabular}{lrrrrrrrr}
\toprule
$k$ & Min k & Edit & Pair & Aux@1 & Field@1 & Exact & Issue & Ret@5 \\
\midrule
1 & 1 & 0.25 & 0.21 & 0.71 & 0.98 & 0.39 & 1.00 & 1.00 \\
2 & 2 & 0.52 & 0.04 & 0.06 & 0.35 & 0.01 & 1.00 & 1.00 \\
3 & 3 & 0.53 & 0.03 & 0.05 & 0.33 & 0.00 & 1.00 & 1.00 \\
5 & 5 & 0.54 & 0.02 & 0.04 & 0.24 & 0.00 & 1.00 & 1.00 \\
8 & 8 & 0.54 & 0.02 & 0.04 & 0.17 & 0.00 & 1.00 & 1.00 \\
12 & 12 & 0.54 & 0.02 & 0.04 & 0.12 & 0.00 & 1.00 & 1.00 \\
20 & 20 & 0.54 & 0.02 & 0.04 & 0.10 & 0.00 & 1.00 & 1.00 \\
\bottomrule
\end{tabular}
\caption{\method{} sensitivity to the target anonymity set size. The main auxiliary matcher drops sharply by $k=2$, while the field-weighted stress attacker continues improving at larger $k$.}
\label{tab:sensitivity}
\end{table}

\subsection{Attack Robustness, RAG Exposure, and API Audit}

Across word, character, hybrid TF-IDF, and field-weighted auxiliary matchers, direct redaction ranges from 0.51 to 0.98 top-1 and the document-local proxy ranges from 0.48 to 0.98.
\method{} remains lower, from 0.03 to 0.24, but the field-weighted result exposes residual structured-context risk.
The residual analysis finds four \method{} top-1 residual matches under the main word attacker; none retain exact quasi-identifier fields, and the median score margin is only 0.00.
This means that \method{} does not eliminate all matching signal, but the remaining matches are weak and mostly coarse.

The RAG-style exposure diagnostic shows why linkage evaluation matters for LLM-ready corpora.
For high-linkage T3 personas, exact profile queries retrieve a target document for every direct-redacted, Presidio-redacted, and document-local corpus; \method{} reduces T3 Hit@5 to 0.31 and Multi@10 to 0.00, while aggressive redaction reduces Hit@5 to 0.03 at much higher utility cost.
Deterministic generated profile-like queries preserve the ordering: verbose-query Hit@5 is 0.38 for direct redaction and the document-local proxy, 0.26 for \method{}, and 0.03 for aggressive redaction; short role-region queries separate direct redaction from \method{} at 0.35 versus 0.02.
A deterministic top-5 RAG context-recovery scan on T3 profiles shows that direct redaction exposes 9.78 exact quasi-identifier fields on average and the document-local proxy exposes 7.84 exact / 9.47 coarse fields, while \method{} exposes 0.00 exact / 0.44 coarse fields.
This diagnostic is not a full RAG attack, but it shows that transformed corpora can remain searchable by auxiliary profiles.

Finally, GPT-5.5 stress audits use synthetic data only and capped calls.
Auxiliary matching preserves the ordering: direct, Presidio, and document-local transformations score 1.00, 0.98, and 1.00 top-1, while \method{} drops to 0.44 with uncertainty 0.92 and aggressive redaction scores 0.33.
A 24-person GPT-5.5 document-local anonymization baseline remains fully matchable at 1.00 top-1.
An additional 24-case GPT-5.5 evidence-extraction audit finds that direct-redaction successful matches cite location in 1.00 and role in 0.75 of cases, while \method{} residual matches cite role, location, and institution at 0.00 and are uncertain in 0.88.
A completed 12-person T3 RAG-generation audit labels direct/Presidio/document-local contexts as same-person at 1.00, versus 0.42 uncertain for \method{} and 0.00 for aggressive redaction.
Because API model behavior can change, these are stress audits rather than the main statistical evidence.

\begin{table}[t]
\centering
\small
\setlength{\tabcolsep}{3pt}
\caption{GPT-5.5 qualitative evidence extraction over 24 selected synthetic cases. Entries are case-level rates for signal families in the extracted explanations.}
\label{tab:gpt55_evidence_signals}
\begin{tabular}{lrrrrrr}
\toprule
Case & n & Role & Loc. & Inst. & High-spec. & Unc. \\
\midrule
Direct success & 8 & 0.75 & 1.00 & 0.50 & 0.47 & 0.00 \\
LG residual & 8 & 0.00 & 0.00 & 0.00 & 0.07 & 0.88 \\
Agg failure & 8 & 0.00 & 0.00 & 0.00 & 0.18 & 0.75 \\
\bottomrule
\end{tabular}
\end{table}

\section{Related Work}

Classic work on $k$-anonymity formalized how released records can remain re-identifiable through combinations of quasi-identifiers and motivated generalization and suppression \citep{sweeney2002}.
Subsequent privacy models such as $\ell$-diversity and $t$-closeness highlighted that anonymity-set size alone is insufficient when sensitive attributes remain homogeneous or distributionally revealing \citep{machanavajjhala2007ldiversity,li2007tcloseness}.
Sparse-data attacks showed that high-dimensional records can be de-anonymized with limited auxiliary knowledge \citep{narayanan2008}, and later work examined how practical de-identification defenses can fail under background knowledge and downcoding attacks \citep{cohen2022attacks}.

Text anonymization benchmarks such as TAB evaluate anonymization decisions and utility in legal text \citep{pilan2022tab}.
RAT-Bench recently reframes text anonymization around re-identification risk from both direct and indirect identifiers \citep{krco2026ratbench}.
Surveys of text anonymization and anonymization metrics emphasize that privacy, utility, legal requirements, and user-centered expectations are often measured inconsistently \citep{deusser2025textanon,ren2025metrics}.
Presidio and OpenAI Privacy Filter illustrate the practical importance of PII detection tools, but such tools are not by themselves complete anonymization guarantees \citep{presidio,openai2026privacyfilter}.

Recent work on LLM privacy shows that language models can infer personal attributes from text beyond direct memorization \citep{staab2024beyond}, and that large models can memorize and reveal training data under extraction attacks \citep{carlini2021extracting,nasr2023scalable}.
RAG privacy work shows that retrieval databases can leak private records through retrieval and generation channels \citep{zeng2024ragprivacy}.
Autoregressive infilling attacks on de-identified documents further show that background knowledge can help recover masked spans \citep{charpentier2025infilling}.
Our contribution is complementary: we focus on cross-document linkage after document-level transformations, where repeated corpus-level quasi-identifiers can survive otherwise reasonable LLM-ready anonymization workflows.

\section{Limitations and Ethics}

The benchmark is synthetic and template-controlled.
This is necessary for safe ground truth and reproducibility, but it limits external validity.
Real sensitive corpora may contain richer writing styles, temporal dependencies, institutional terminology, missing fields, contradictions, OCR errors, and adversary knowledge.
Our local attacks are lightweight, and
\method{} is a heuristic generalization method.
The field-weighted attacker shows that residual generalized context can still be informative.
Utility is measured with issue labels, retrieval, and lightweight fact checks rather than full downstream deployment tasks.

The benchmark is intended for defensive auditing.
We do not collect, process, or attempt to re-identify real people.
All examples are synthetic and generated from a controlled schema.
Any extension to real data would require authorization, data minimization, access controls, logging, and ethical review.
Practitioners should not treat successful performance on this benchmark as proof that a real corpus is safe to release or train on.

\section{Conclusion}

Document-local PII removal is not enough for LLM-ready sensitive corpora.
In a multi-document setting, repeated quasi-identifiers can support linkage, auxiliary-profile matching, and profile-query retrieval even after direct identifiers are removed.
Our synthetic benchmark suggests that privacy evaluation should include corpus-level linkage resistance, not only span-level PII recall.
It also suggests that simple corpus-aware generalization can improve the privacy--utility frontier relative to document-local anonymization and blanket redaction.
The broader lesson is that data transformation for foundation models should ask what a capable attacker can still infer across the corpus, not only what identifying spans were removed from each document.

\bibliographystyle{colm2026_conference}
\begin{thebibliography}{16}

\bibitem[Carlini et~al.(2021)]{carlini2021extracting}
Nicholas Carlini, Florian Tram{\`e}r, Eric Wallace, Matthew Jagielski, Ariel Herbert-Voss, Katherine Lee, Adam Roberts, Tom Brown, Dawn Song, {\'U}lfar Erlingsson, Alina Oprea, and Colin Raffel.
\newblock Extracting training data from large language models.
\newblock In \emph{USENIX Security Symposium}, 2021.

\bibitem[Charpentier and Lison(2025)]{charpentier2025infilling}
Lucas Charpentier and Pierre Lison.
\newblock Re-identification of de-identified documents with autoregressive infilling.
\newblock \emph{arXiv:2503.03438}, 2025.

\bibitem[Cohen(2022)]{cohen2022attacks}
Aloni Cohen.
\newblock Attacks on deidentification's defenses.
\newblock In \emph{USENIX Security Symposium}, 2022.

\bibitem[Deu{\ss}er et~al.(2025)]{deusser2025textanon}
Theresa Deu{\ss}er, Andreas Kirsch, Florian Matthes, and Simon Mayer.
\newblock SoK: A survey of text anonymization techniques.
\newblock \emph{arXiv:2508.15433}, 2025.

\bibitem[Kr{\v{c}}o et~al.(2026)]{krco2026ratbench}
Nata{\v{s}}a Kr{\v{c}}o, Zexi Yao, Matthieu Meeus, and Yves-Alexandre de Montjoye.
\newblock RAT-Bench: A comprehensive benchmark for text anonymization.
\newblock \emph{arXiv:2602.12806}, 2026.

\bibitem[Li et~al.(2007)]{li2007tcloseness}
Ninghui Li, Tiancheng Li, and Suresh Venkatasubramanian.
\newblock $t$-closeness: Privacy beyond $k$-anonymity and $\ell$-diversity.
\newblock In \emph{IEEE International Conference on Data Engineering}, 2007.

\bibitem[Machanavajjhala et~al.(2007)]{machanavajjhala2007ldiversity}
Ashwin Machanavajjhala, Daniel Kifer, Johannes Gehrke, and Muthuramakrishnan Venkitasubramaniam.
\newblock $\ell$-diversity: Privacy beyond $k$-anonymity.
\newblock \emph{ACM Transactions on Knowledge Discovery from Data}, 2007.

\bibitem[Microsoft(2026)]{presidio}
Microsoft.
\newblock Presidio: Data protection and de-identification SDK.
\newblock \url{https://microsoft.github.io/presidio/}, 2026.

\bibitem[Narayanan and Shmatikov(2008)]{narayanan2008}
Arvind Narayanan and Vitaly Shmatikov.
\newblock Robust de-anonymization of large sparse datasets.
\newblock In \emph{IEEE Symposium on Security and Privacy}, 2008.

\bibitem[Nasr et~al.(2023)]{nasr2023scalable}
Milad Nasr, Nicholas Carlini, Jonathan Hayase, Matthew Jagielski, A. Feder Cooper, Daphne Ippolito, Christopher A. Choquette-Choo, Eric Wallace, Florian Tram{\`e}r, and Katherine Lee.
\newblock Scalable extraction of training data from production language models.
\newblock \emph{arXiv:2311.17035}, 2023.

\bibitem[OpenAI(2026)]{openai2026privacyfilter}
OpenAI.
\newblock Introducing OpenAI Privacy Filter.
\newblock \url{https://openai.com/index/introducing-openai-privacy-filter/}, 2026.

\bibitem[Pil{\'a}n et~al.(2022)]{pilan2022tab}
Ildik{\'o} Pil{\'a}n, Pierre Lison, Lilja {\O}vrelid, Anthi Papadopoulou, David S{\'a}nchez, and Montserrat Batet.
\newblock The Text Anonymization Benchmark (TAB): A dedicated corpus and evaluation framework for text anonymization.
\newblock \emph{Computational Linguistics}, 2022.

\bibitem[Ren et~al.(2025)]{ren2025metrics}
Hao Ren, Yi Zhang, and Yang Li.
\newblock A survey of text anonymization metrics.
\newblock \emph{arXiv:2511.06282}, 2025.

\bibitem[Staab et~al.(2024)]{staab2024beyond}
Robin Staab, Mark Vero, Mislav Balunovi{\'c}, and Martin Vechev.
\newblock Beyond memorization: Violating privacy via inference with large language models.
\newblock In \emph{International Conference on Learning Representations}, 2024.

\bibitem[Sweeney(2002)]{sweeney2002}
Latanya Sweeney.
\newblock Achieving $k$-anonymity privacy protection using generalization and suppression.
\newblock \emph{International Journal of Uncertainty, Fuzziness and Knowledge-Based Systems}, 2002.

\bibitem[Zeng et~al.(2024)]{zeng2024ragprivacy}
Shenglai Zeng, Jiankun Zhang, Pengfei He, Yiding Liu, Yiming Liu, Hanxu Hu, Lianhui Qin, Yue Xing, Mengdi Wang, and Jiliang Tang.
\newblock The good and the bad: Exploring privacy issues in retrieval-augmented generation.
\newblock \emph{arXiv:2402.16893}, 2024.

\end{thebibliography}

\end{document}
